\section{Introduction}
\label{sec:introduction}

A customer asks to return a customized item within the 30-day window. The agent consults the store policy, which states that ``items may be returned within 30 days.'' The word ``items'' scopes over all product categories, including custom orders that are normally final sale. The agent grants a full refund and arranges free return shipping. One word of scopal ambiguity costs the retailer \$200.

LLM-based agents increasingly operate under natural-language policy specifications that govern their permitted actions. These specifications inevitably contain ambiguity: a property that has been studied for decades in requirements engineering (RE) and computational linguistics. \citet{berry_kamsties_2004} distinguish between \emph{nocuous} (i.e., harmful) ambiguity, where different stakeholders arrive at different interpretations, and \emph{innocuous} ambiguity, where shared organizational context ensures convergent interpretation. In human RE workflows, most ambiguity is innocuous because developers share domain knowledge, can consult colleagues, and iteratively clarify requirements.

LLM agents lack these disambiguation channels. They cannot consult colleagues, rarely escalate to users, and must resolve ambiguity in real time from the text alone. This raises a question that has not been empirically addressed: does policy ambiguity cause safety violations in LLM agents, and if so, does the \emph{type} of ambiguity predict the severity of the violation? Prior work classifies ambiguity types in policies \citep{drip_r_2026,li_ambiguity_taxonomy_2024} or measures policy understanding quality \citep{policybank_2026}, but no existing study provides a causal, type-granularity mapping from ambiguity to safety violations.

\begin{figure*}[t]
\centering
\includegraphics[width=\textwidth]{figures/fig1_overview.pdf}
\caption{Experimental overview. We construct 300 matched clause pairs across a six-type, two-layer ambiguity taxonomy and inject them into $\tau^2$-bench policy documents. Five models interact with the modified policies across 2,997 episodes. A cross-judge evaluation pipeline classifies violation severity on a four-tier scale.}
\label{fig:overview}
\end{figure*}

We address this gap with a matched-pair causal study on $\tau^2$-bench \citep{tau2_bench_2025}, a dual-control conversational agent benchmark. Using a six-type taxonomy organized into two layers (linguistic and specification, \S\ref{sec:taxonomy}), we construct 300 ambiguous/unambiguous clause pairs and evaluate five architecturally diverse models across 2,997 episodes.

Our results show that all policy ambiguity is nocuous for LLM agents ($+34.0$ pp violation rate increase, OR $= 15.1$), but that type primarily determines severity, with a smaller but significant model$\times$type interaction (2.9\% of variance). Incompleteness is consistently high (46--62\% violation rate across all models), while authorization scope shows the widest but least reliably measured cross-model range (52\,pp; $\kappa = 0.385$, \S Limitations), suggesting that safety training may differentially mitigate specific ambiguity types. Agents exhibit disambiguation deficit in over 89\% of cases, with fewer than 11\% showing any disambiguation awareness.

Our contributions are: (1)~the first causal quantification of policy ambiguity's safety impact via matched-pair design with perturbation control (\S\ref{sec:method}), observed across five primary models and confirmed in two additional replication models; (2)~incompleteness as a consistent upper-bound risk type across all models and margins, with type identity explaining 22$\times$ more variance than model identity (\S\ref{sec:results}); and (3)~a failure mode taxonomy with the empirical phenomenon of \emph{Disambiguation Deficit} ($\leq$11\% awareness; \S\ref{sec:analysis}).
